\subsubsection{Date de intrare}
Conform specificațiilor organigramei, parametrii sistemului sunt:
\begin{itemize}
    \item Număr de stări ($N_{stari}$): 22
    \item Număr variabile de intrare ($N_{in}$): 9
    \item Număr variabile de ieșire ($N_{out}$): 5
\end{itemize}

\subsubsection{Dimensionarea câmpurilor}

\begin{enumerate}
    \item \textbf{Câmpul de adresare ($n_{adr}$):}
    Numărul de biți necesari pentru a adresa unic cele 22 de stări se calculează logaritmic:
    \[
    n_{adr} = \log_2(N_{stari}) = \log_2(22)
    \]
    Deoarece $2^4 = 16 < 22 \leq 32 = 2^5$:
    \[
    n_{adr\_necesar} = 5 \text{ biți}
    \]
  Vom alege memoria de tip \textbf{64 x 4} pentru implementarea microinstrucțiunilor așa încât:
    \[
    n_{adr} = \log_2(64) = 6 \text{ biți}
    \]

    \item \textbf{Câmpul condițiilor de intrare ($n_{ci}$):}
    Pentru a selecta una dintre cele 9 variabile de intrare (condiții de test), calculăm:
    \[
    n_{ci} = \log_2(N_{in}) = \log_2(9)
    \]
    Deoarece $2^3 = 8 < 9 \leq 16 = 2^4$:
    \[
    n_{ci} = 4 \text{ biți}
    \]

    Astfel va fi folosit un \textbf{MUX 16:1} pentru schema unității de comandă.

    \item \textbf{Câmpul de ieșire ($n_{out}$):}
    În formatul variabil, pentru instrucțiunile de tip 0 (comandă), biții sunt alocați direct variabilelor de ieșire.
    \[
    n_{out} = 5 \text{ biți}
    \]
\end{enumerate}

\subsubsection{Aplicarea calculului final}

Lungimea microinstrucțiunii este determinată de maximul dintre lungimea necesară pentru instrucțiunile de salt (Tip 1) și cele de comandă (Tip 0), plus bitul de selecție a formatului:

\[
L_{\mu i} = 1 + \max(\underbrace{n_{ci} + n_{adr}}_{\text{Tip 1}}, \underbrace{n_{out}}_{\text{Tip 0}})
\]

Înlocuind valorile calculate:

\[
L_{\mu i} = 1 + \max(4 + 6, 5)
\]
\[
L_{\mu i} = 1 + \max(9, 5)
\]
\[
L_{\mu i} = 1 + 9 = 11 \text{ biți}
\]

\textbf{Concluzie:} Lungimea microinstrucțiunii în format variabil este de \textbf{11 biți}, structurată astfel:
\begin{itemize}
    \item \textbf{Bit 10:} Selector format (1 bit).
    \item \textbf{Biții 9-0 (Tip 1):} 4 biți pentru condiționale + 6 biți pentru adresele stărilor (primul bit pentru adresele stărilor va fi mereu 0, deoarece sunt doar 5 biți de adresă necesari).
    \item \textbf{Biții 9-0 (Tip 0):} 5 biți utilizați pentru Ieșiri (restul de 5 biți fiind neutilizați).
\end{itemize}

  Vom concatena \textbf{3 circuite} de memorie \textbf{64 x 4} pentru implementarea dispozitivului.
